% BANKS-SOURCE: pdf=231 print=221 kind=prose
that the strength of the potential is exponentially small for small $g$. This is
quite different from the phase where there is no expectation value of the Higgs
field. There the strength of the Coulomb potential is of order $g^2$.

In the interpretation of our partition function and Wilson loop in terms of the
Aharonov--Bohm phase in a superconductor, the failure of screening in the
presence of instantons is an indication that the Meissner effect disappears at
finite temperature. Two-dimensional superconductors have finite-energy
excitations that are ``Abrikosov flux dots.'' At finite temperature the system is
a dilute gas of dots and anti-dots. An external magnetic field penetrates the
system by preferentially aligning the dots in the region where the field exists.

% BANKS-SOURCE: pdf=231 print=221 kind=section id=10.5
\subsection{Monopole instantons in three-dimensional Higgs models}

% BANKS-SOURCE: pdf=231 print=221 kind=prose
In three Euclidean dimensions, topological charges are classified by
$\Pi_2(M_V)$, the mappings of the 2-sphere into the space of vacuum field
configurations. The space of gauge field vacua is always the gauge group, and
$\Pi_2(G)$ vanishes for any Lie group. Thus, we will have instantons only for
models that involve scalar fields, and only if the space of minima of the scalar
potential has non-vanishing $\Pi_2$. The angular-momentum part of the kinetic
term has the form

% BANKS-SOURCE: pdf=231 print=221 kind=display
\[
  \int\dd r\,\Phi^*L^2\Phi,
\]

so angular dependence leads to a linearly divergent action. We conclude, as in
two dimensions, that finite-action instantons will occur only if the space of
scalar vacua has a gauge equivalence on it. That is, locally it takes the form
$G/H\times X$, where $G$ is the gauge group and $H$ the subgroup which
preserves a point in the vacuum manifold. In physics language, the gauge group
$G$ is in the Higgs phase, with $H$ the unbroken subgroup. In our discussion
below of magnetic monopoles as solitons, we will show that
$\Pi_2(G/H)=\Pi_1(H)$. The most interesting case will be that in which $H$ is a
product of $U(1)$ groups, and the simplest example is a single $U(1)$ and
$G=SU(2)$. In this case, $G/H$ is just $S^2$ and the topological charge is just
the winding of the sphere on itself. The map of charge 1 is just the identity
map, while that of charge $-1$ is the orientation-reversing map of the sphere on
itself, $n^a\to-n^a$, in the representation of $S^2$ as the space of unit
3-vectors.

Thus, the simplest three-dimensional model with finite-action instantons is the
Georgi--Glashow model of $SU(2)$ gauge theory broken to $U(1)$ by a Higgs field
in the three-dimensional adjoint representation. This is the Euclidean
instanton sector whose static soliton interpretation is the $D=4$ Lorentzian
monopole theory. The action is

% BANKS-SOURCE: pdf=231 print=221 kind=display
% CONVENTION-SCOPE: fixed D=4 reason=three-dimensional Euclidean monopole instanton and D=4 Lorentzian soliton
\[
  S=\int\dd^3\mathbf x\,\frac{1}{4g^2}
  \left[(F^a{}_{\mu\nu})^2+\frac12(D^{ab}{}_{\mu}\phi_b)^2
  +V(\phi^a\phi^a)\right].
\]

% BANKS-SOURCE: pdf=232 print=222 kind=prose
We count the mass dimension of both the gauge potential and the scalar field as
1, in which case $g^2$ also has mass dimension 1. These are not the scaling
dimensions of the fields in a renormalization-group analysis. Restricting our
attention to renormalizable potentials, $V$ is a polynomial of order $\leq6$
whose minimum is at non-zero $\phi^a$. In order to get explicit solutions, we
will restrict ourselves to the case $V=0$, a situation that is natural if there
is enough supersymmetry. None of the qualitative results we obtain will depend
on this restriction. The vacuum expectation value of $|\phi^a|$ is denoted $\nu$.

Our analysis begins from the remark that

% BANKS-SOURCE: pdf=232 print=222 kind=display
% CONVENTION-SCOPE: fixed D=4 reason=three-dimensional Euclidean monopole instanton and D=4 Lorentzian soliton
\[
  \int\dd^3\mathbf x\,\frac14
  \left[F^a{}_{\mu\nu}\pm\epsilon_{\mu\nu\lambda}(D_\lambda\phi)^a\right]^2
  \geq0.
\]

Thus, if the potential is zero,

% BANKS-SOURCE: pdf=232 print=222 kind=display
% CONVENTION-SCOPE: fixed D=4 reason=three-dimensional Euclidean monopole instanton and D=4 Lorentzian soliton
\[
  g^2S\pm\frac12\int\dd^3\mathbf x\,
  F^a{}_{\mu\nu}\epsilon_{\mu\nu\lambda}(D_\lambda\phi)^a\geq0.
\]

It follows that

% BANKS-SOURCE: pdf=232 print=222 kind=display
% CONVENTION-SCOPE: fixed D=4 reason=three-dimensional Euclidean monopole instanton and D=4 Lorentzian soliton
\[
  S\geq|M|\frac{\nu}{g^2},
\]

where the magnetic charge $M$ is defined by

% BANKS-SOURCE: pdf=232 print=222 kind=display
% CONVENTION-SCOPE: fixed D=4 reason=three-dimensional Euclidean monopole instanton and D=4 Lorentzian soliton
\[
  M\nu=\frac12\int\dd^3\mathbf x\,
  F^a{}_{\mu\nu}\epsilon_{\mu\nu\lambda}(D_\lambda\phi)^a.
\]

$\nu$ is the vacuum expectation value of the dimension-1 Higgs field. Now we
can use the Leibniz rule for covariant derivatives, the Bianchi identity for the
gauge field strength, and the fact that covariant derivatives of singlets are
ordinary derivatives, to show that

% BANKS-SOURCE: pdf=232 print=222 kind=display
\[
  M=\frac1\nu\int_{S^2}\phi^aF^a,
\]

where $F^a=\tfrac12F^a{}_{\mu\nu}\dd x^\mu\dd x^\nu$ is the field-strength
2-form.
\BanksImplicitHook{I-CH10-012}

The name magnetic charge comes from the fact that we have an unbroken $U(1)$
gauge theory. If we interpret our solutions as static solutions of the
$D=4$, $(3+1)$-dimensional version of the theory, then $M$ measures the magnetic flux
of the unbroken gauge field through the sphere at infinity. These solutions are
called 't Hooft--Polyakov monopoles.

The inequality is saturated by solutions of the first-order equations

% BANKS-SOURCE: pdf=232 print=222 kind=display
\[
  \frac12\epsilon_{\lambda\mu\nu}F^a{}_{\mu\nu}
  =\pm(D_\lambda\phi^a).
\]

% BANKS-SOURCE: pdf=232 print=222 kind=prose
In the supersymmetric context which justifies setting $V=0$ in the quantum
theory, these are the Bogomol'nyi--Prasad--Sommerfield (BPS) equations which say
that the

% BANKS-SOURCE: pdf=233 print=223 kind=prose
solution preserves half the supercharges.\footnote[8]{This is most easily
understood in the soliton interpretation of our solutions in $3+1$ dimensions.
The extended supersymmetry algebra
\[
  [Q^i{}_\alpha,\bar Q^j{}_\beta]_+
  =(\gamma^\mu)_{\alpha\beta}P_\mu\delta^{ij}
  +\epsilon^{ij}M
\]
has degenerate representations when the matrix
$\gamma^\mu P_\mu+\epsilon M$ has zero eigenvalues. This is a condition on the
mass of particles, which for solitons is the static energy. The energy becomes
Euclidean action in the three-dimensional instanton interpretation of the
solutions.} This accounts for the first-order nature of the equations: SUSY
variations are of first order in derivatives.

In the theory with no potential there are $n$ monopole solutions. These arise
because the massless scalar field which parametrizes the radius of $\phi^a$ gives
rise to attractive long-range forces that cancel out the magnetic repulsion of
the monopoles. For our purposes we will be more interested in the non-BPS
solutions consisting of equal numbers of monopoles and anti-monopoles. The
properties of these can be completely understood in terms of the BPS solution
with unit magnetic charge. It is natural to make a spherically symmetric ansatz,

% BANKS-SOURCE: pdf=233 print=223 kind=display
\[
  \phi^a=\frac{x^a}{r}f(r),
  \qquad
  A^a{}_{\mu}=\epsilon_{\mu a\nu}\frac{x^\nu}{r}a(r).
\]

Note that $f$ and $a$ have mass dimension 1.

We have

% BANKS-SOURCE: pdf=233 print=223 kind=display
\[
  (D_\lambda\phi)^a=P^a{}_{\lambda}
  \left(\frac1r-a\right)f+\frac{x^ax_\lambda}{r^2}f',
\]

where $P_{ab}$ is the projector orthogonal to $x_a/r$.

The field strength is given by

% BANKS-SOURCE: pdf=233 print=223 kind=display
\[
  2B^a{}_\lambda\equiv\epsilon_{\mu\nu\lambda}F^a{}_{\mu\nu}
  =2P_{\lambda a}\left(-\frac a r-a'\right)
  +2\frac{x_ax_\lambda}{r^2}
  \left(a^2-\frac{2a}{r}\right).
\]

The Bogomol'nyi equation

% BANKS-SOURCE: pdf=233 print=223 kind=display
\[
  B^a{}_{\lambda}=\pm(D_\lambda\phi^a)
\]

leads to

% BANKS-SOURCE: pdf=233 print=223 kind=display
\[
  \pm\left(\frac1r-a\right)f=-\left(\frac a r+a'\right)
\]

and

% BANKS-SOURCE: pdf=233 print=223 kind=display
\[
  \pm f'=a^2-\frac{2a}{r}.
\]

If we introduce

% BANKS-SOURCE: pdf=233 print=223 kind=display
\[
  q\equiv(1-ra),
\]

then the first equation reads $f=q'/q$, and the equations are solved by

% BANKS-SOURCE: pdf=233 print=223 kind=display
\[
  q=\frac{r\nu}{\sinh(r\nu)},
\]

where $\nu$ is the vacuum expectation value of $|\phi^a|$.

% BANKS-SOURCE: pdf=234 print=224 kind=prose
We will do the calculation of the effects of instantons in the low-energy
effective theory below the mass of the massive gauge bosons and scalars. In a
generic theory with monopole instantons, this low-energy theory consists of a
three-dimensional $U(1)$ gauge theory, coupled to a dilute gas of point-like
instantons. If we think of the vector dual to $F_{ij}$ as an *electric field* in
the statistical mechanics of three-dimensional electrostatics, then our problem
is equivalent to the statistical mechanics of a Coulomb gas. In that language,
the phenomenon we are about to expose is called *Debye screening*. For a given
configuration of monopoles and anti-monopoles, the field strength is given by

% BANKS-SOURCE: pdf=234 print=224 kind=display
\[
  F^{\rm mon}_{\mu\nu}=\epsilon_{\mu\nu\lambda}\partial_\lambda\phi_0,
\]

where

% BANKS-SOURCE: pdf=234 print=224 kind=display
\[
  \nabla^2\phi_0=4\pi\sum
  \left[\delta^3(\mathbf x-\mathbf x_i)-\delta^3(\mathbf x-\mathbf y_i)\right].
\]

$\mathbf x_i$ and $\mathbf y_i$ are the positions of the monopoles and anti-monopoles,
respectively. We have to sum over all monopole and anti-monopole numbers, with
weight

% BANKS-SOURCE: pdf=234 print=224 kind=display
\[
  \frac{1}{N_+!N_-!}
  \left[D^{-1/2}
  \left(\frac{\mathcal N}{2\pi(gR_c)^2}\right)^{3/2}
  \ee^{-\frac{S_0}{(gR_c)^2}}\right]^{N_++N_-},
\]

and integrate over all the positions. $\mathcal N$ is the integral of the
square of the translational zero modes (summed over all field components). This
quantity, as well as $D$ and $S_0$, are the contributions of the microscopic
theory to the instanton amplitudes. The $R_c$ in this formula is $1/\nu$. In
principle there are renormalization corrections to this effective field theory
formula, but if $gR_c$ is small, then they are small, and we can neglect them.

Different underlying theories, with different scalar potentials, or other
massive fields, will change these parameters, but will not otherwise affect the
infrared dynamics we study. The parameters should be tuned to absorb infinite
effects in the low-energy theory, such as the Coulomb self-energies of the
monopoles, and restore the finite values of these self-energies in the underlying
theory. Additional massless fields, such as those guaranteed in a supersymmetric
theory, can change the IR dynamics. Typically the change has to do with the
particular Green functions to which the instantons contribute, rather than a
drastic change in the instanton statistical mechanics.

The fluctuating $U(1)$ gauge dynamics can be represented by a vector potential in
a particular gauge. It is more convenient, however, to represent it by a
fluctuating field strength $\mathcal F_{\mu\nu}$. This is achieved by replacing
the Euclidean Maxwell action by

% BANKS-SOURCE: pdf=234 print=224 kind=display
\[
\frac{g^2}{4}\mathcal F_{\mu\nu}^{2}
  +2\ii\mathcal F_{\mu\nu}(\partial_\nu A_\mu)
  +\text{(gauge fixing)}
\]

(note that $\mathcal F$ has dimension 1). Integrating over $\mathcal F_{\mu\nu}$
restores the Maxwell form, while integrating over $A_\mu$ gives us a functional
delta function setting

% BANKS-SOURCE: pdf=234 print=224 kind=display
\[
  \partial_\nu\mathcal F_{\mu\nu}=J_\mu.
\]

% BANKS-SOURCE: pdf=235 print=225 kind=prose
The current is zero in the denominator of the functional integral formula, and
represents the electrically charged particles (in the sense of $(2+1)$-dimensional
electrodynamics) coupled to the gauge field. We will take it to be the current
of a single Wilson loop.

The constraint is solved by setting

% BANKS-SOURCE: pdf=235 print=225 kind=display
\[
  \mathcal F_{\mu\nu}=\mathcal F^0_{\mu\nu}
  +\epsilon_{\mu\nu\lambda}\partial_\lambda\phi,
\]

where $\mathcal F^0$ is any special solution of the constraint equation. $\phi$
is a dimensionless scalar field. Different choices of $\mathcal F^0$ are
absorbed into the $\phi$ functional integral. Later, we will make a
semi-classical approximation for $\phi$, and it is convenient to choose
${\mathcal F}^0$ to simplify the semi-classical analysis. We will choose it to be
zero in the denominator and, in the numerator, a surface delta function on the
minimal-area surface bounded by the loop.

The utility of the dual formulation is that it is easy to include the monopole
instantons as sources of the fluctuating field $\phi$. If we add the term

% BANKS-SOURCE: pdf=235 print=225 kind=display
\[
  \ii\sum_i[\phi(\mathbf x_i)-\phi(\mathbf y_i)]
\]

to the denominator functional integral, integration over $\phi$ gives rise to
the repulsive and attractive long-range Coulomb forces among monopoles and
anti-monopoles. It also leads to infinite self-energies, which are absorbed into
$S_0$. Note that the factor of $\ii$ in the action is required in order to get
the right sign for the Coulomb energies and make like charges repel and opposite
charges attract.

We can now sum over the monopoles for fixed $\phi$ and obtain a correction to
the effective action

% BANKS-SOURCE: pdf=235 print=225 kind=display
\[
  \delta S=2D^{-1/2}
  \left(\frac{\mathcal N}{2\pi(gR_c)^2}\right)^{3/2}
  \ee^{-\frac{S_0}{(gR_c)^2}}\cos\phi,
\]

a result first obtained by Polyakov (though the equivalent physics was done long
ago by Debye). The most significant aspect of this result is that the $\phi$
field has obtained an exponentially small mass.

Now consider trying to evaluate the functional integral over $\phi$ by finding a
classical solution that minimizes the action. In the denominator, the appropriate
solution is just $\phi=0$. In the numerator, if we ignore the cosine term, we
could try to set

% BANKS-SOURCE: pdf=235 print=225 kind=display
\[
  \partial_\lambda\phi=\frac12\epsilon_{\mu\nu\lambda}
  \mathcal F^0_{\mu\nu}.
\]

This cannot work everywhere, because of the current source on the Wilson loop.
However, it fails only in the vicinity of the loop and we get an action
proportional to the circumference of the loop. The large-distance behavior of the
loop is in this case determined by the quadratic fluctuations of $\phi$ around
the classical solution, which give rise to the logarithmically rising Coulomb
potential of $(2+1)$-dimensional electrodynamics.

For simplicity, think about a loop in the $x,y$ plane, with $z=0$. The solution
discussed above has $\phi=0$ for negative $z$. It jumps to
$\phi=\tfrac12\mathcal F^0_{xy}$, for positive $z$, whenever $(x,y)$ is in the
interior of the loop. In general, this will lead to an infinite contribution
from the cosine term, of the form $A[\infty]$, where $A$ is the area of the loop.
We have to let $\phi$ go back to zero to avoid the infinity, but this leaves over
a finite contribution proportional to $A$ for the minimal-action configuration.
The proportionality constant is exponentially small.

Note, however, that, when the required jump in $\phi$ is exactly an integral
number of periods of the cosine, we can eliminate the area term entirely. By
arithmetic that

% BANKS-SOURCE: pdf=236 print=226 kind=prose
we will do carefully when we study monopoles as solitons, we find that the area
law disappears precisely for Wilson loops corresponding to the charges carried by
massive $W$ bosons. Thus we get a linear confining potential only between
charges corresponding to half-integral spin representations of $SU(2)$.

The answer to the question of what classical configurations these monopole
instantons are tunneling between is quite subtle. To find it we compactify the
two space dimensions on a rectangular torus of radius $R$. The gauge-invariant
magnetic field $\phi^aF^a{}_{12}$ has configurations with non-vanishing quantized
flux wrapping the torus,

% BANKS-SOURCE: pdf=236 print=226 kind=display
\[
  \int_{T^2}\phi^aF^a=N,
\]

and these configurations have energy $\sim1/R^2$. Thus, in the large-$R$ limit,
the system has degenerate states with different values of magnetic flux, and the
monopoles are the instantons which tunnel between these states.

% BANKS-SOURCE: pdf=236 print=226 kind=section id=10.6
\subsection{Yang--Mills instantons}

\setcounter{footnote}{8}
